% !TEX root = main.tex
% ---------------------------------------------------------------------------
% Introduction --- v2 (2026-08-26). Restructured from 6 paragraphs to 3 at the
% user's request, and shortened (~1350 -> ~780 words).
%   P1  background + the problem/challenge
%   P2  prior work / standard solutions, and what each lacks
%   P3  what we propose + contributions
%
% Positioning notes (do not lose these):
%  * \citep{ikemura2025stressrl} is OUR OWN PRIOR ITERATION (rejected, arXiv
%    public). P2 must state the delta from it explicitly and early, or the paper
%    reads as incremental. Delta: stress as a GRASP-PLANNING objective at
%    demonstration time (IL) vs a stress-penalised RL REWARD; and the problem
%    addressed is COVERAGE across object variation, which RL per-task
%    optimisation does not address. If the venue is double-blind, cite in third
%    person.
%  * \citep{pan2020stressmin} already established stress as a grasp metric. We
%    do NOT claim that idea. Our deltas: a POSITION(width)-controlled contact
%    model matching a real parallel jaw, the E-linearity that makes stiffness
%    randomisation ~free, and use as a demonstration generator.
% Deliberately NOT claimed: lead/dwell BC failure modes; "sim ranking inverts";
% measured real gentleness; any tactile baseline comparison.
% ---------------------------------------------------------------------------
\section{Introduction}
\label{sec:intro}

Picking up a mushroom without bruising it is harder than it looks, and damage-free robotic
handling of fragile produce remains an open problem in agricultural and food
automation~\citep{wang2025damageless}. Two quantities have to be right at once. The gripper
must close to a width that is wide enough that the object does not slip and narrow enough that
it does not crush, and for soft produce the margin between the two is small
\textbf{[TODO: measured drop/crush band, A2]}. It must also close in the right place and
orientation: fingers that straddle the stem, or that close across the thin rim of the cap,
fail or bruise at a width that would have been correct elsewhere on the same specimen. Both
the grasp pose and the width are determined by the object's own pose, size, shape and
stiffness, and none of these is directly available to a robot observing the scene with a depth
camera. Pose and size are only partially visible under occlusion, shape varies continuously
within a single produce category, and stiffness is not a geometric property at all. The task
looks like pick-and-place but behaves like a precision problem, in which the quantities that
must be predicted accurately are ones no sensor in the loop reports. We study this on a 7-DoF
arm with a parallel-jaw gripper, a single fixed external depth camera, and no tactile or force
sensing in the control loop, using the lifting of a fresh mushroom as a running example.

Three lines of work address this setting, and each leaves the same thing unresolved.
\emph{Compliant hardware} absorbs the problem mechanically: pneumatic and enveloping grippers
limit contact pressure by construction and reduce bruising~\citep{wang2024softgripper}, but
compliance does not decide where on the object to close, and it is a fixed property of the
end-effector rather than a decision the robot makes per instance. \emph{Contact sensing} gives
the policy access to grip state at execution time, and recent tactile and force-controlled
systems handle delicate items well by this route~\citep{tfgripper,vtam}. This changes what the
policy can observe; it does not change where its supervision comes from. Both are imitation
learners trained on human demonstrations recorded on real hardware, and the division of labour
inside them is telling: the force-regulation policy of \citet{tfgripper} is decoupled from
arm-pose prediction, and the pose policy remains demonstration-learned, so contact sensing
addresses how firmly to close but not where or how to approach. (Deformation-based tactile is
additionally ill-conditioned in the gentle regime, since it needs measurable indentation to
produce a signal, a limitation recognised within the tactile community
itself~\citep{lighttact}.) \emph{Stress-aware methods} target the quantity that actually causes
damage. \citet{pan2020stressmin} rank grasps by the internal stress a held object sustains,
under a point-contact force-closure model. We implemented that metric and found it degenerate
in our setting: with point contacts, two pads on a smooth organic object cannot span the
moment axes of wrench space, so the resistible-wrench hull collapses and the metric returns
essentially zero for \emph{every} candidate grasp on our objects, while a cube clears it
(Section~\ref{sec:method:synth}). A parallel jaw does not make point contact, and the
distinction is not cosmetic. In our own earlier iteration we instead trained a policy with a
stress-penalised reward and transferred it zero-shot~\citep{ikemura2025stressrl}; extending
its evaluation revealed that it did not adapt to changes in object orientation or size, since
reinforcement learning must \emph{explore} that variation rather than be shown it, and is too
sample-inefficient to cover it. Both experiences point the same way. What none of these lines
supplies is coverage of the physical variation that determines the correct grasp in the first
place. That variation cannot be covered by
demonstration: a trajectory is recorded against one specimen in one pose and cannot be
transported to another, so cost scales with the product of the varying dimensions and
multiplies again across categories. For precision-limited tasks the requirement is
super-exponential in the target precision, $\log N \propto 1/(P-c)$
\citep{curseofprecision} \textbf{[TODO: verify claim and venue, C1]}, so low-data real
collection is not merely inconvenient here but structurally inadequate. We observe exactly the
predicted behaviour: a policy trained on 50 real mushroom demonstrations does not degrade
uniformly on hardware but fails in a \emph{coverage-shaped} way, misgrasping at particular
workspace locations and over- or under-squeezing at particular combinations of object size and
location (Figure~\ref{fig:coverage}).

We therefore treat coverage, rather than sensing or reward design, as the binding constraint,
and obtain it from simulation. Simulation buys nothing unless demonstrations can be produced
inside it without a human, which requires an operational definition of gentleness that a
search can optimise; internal stress supplies one, being unobservable on hardware but exact in
a deformable simulator. We score a candidate grasp by the stress it induces under a
\emph{width-controlled, distributed-pad} contact model: a real parallel jaw commands a
displacement over a finite pad, and the resulting force is an outcome. Replacing point
contacts with a pad, and full wrench resistance with a friction margin sufficient to hold the
object against gravity, removes the degeneracy above and yields a score that discriminates
among grasps on exactly the objects where the force-closure formulation cannot. We search the
7-DoF grasp (pose and commanded width) with CMA-ES under that score. Because the
displacement is prescribed, stress is linear in Young's modulus, so a single finite-element
solve per pose covers an entire distribution of stiffnesses and material randomisation becomes
almost free. This is what turns a grasp metric into a coverage instrument: the planner runs
over randomised pose, size, shape and material, and its demonstrations are distilled into a
point-cloud diffusion policy acting from a single depth view, co-trained with a small slice of
real demonstrations that calibrates the residual sensor gap. Our contributions are
(i)~an autonomous stress-minimising demonstrator built on a width-controlled, distributed-pad
contact model, which discriminates among parallel-jaw grasps on organic objects where the
point-contact force-closure formulation is degenerate, and whose stiffness-linearity makes
randomised material coverage tractable (Section~\ref{sec:method:synth})
\textbf{[TODO: confirm the planner comparison is the paired form unaffected by the mesh-scale
bug, C3]}; (ii)~evidence that coverage rather than demonstration count is the binding
constraint, with a policy trained on 50 real demonstrations failing in a coverage-shaped way
while broad simulated coverage co-trained with a small real slice outperforms both
pure-simulation and real-only policies \textbf{[TODO: $n$ and CI, A3; confirm the run
postdates the perception-bias correction, A6]}; and (iii)~real-robot validation across object
poses, sizes and shapes, and across object categories, from a single depth view with no
tactile or force sensing in the loop \textbf{[TODO: generalist run, B2 --- blocking, this is a
stated rejection reason]}. We are explicit about the scope of the word
\emph{gentle}: it is established here under the simulator's stress model and on hardware by
task success, not by a direct measurement of damage \textbf{[TODO: real gentleness proxy, A1]}.
